\documentclass[10pt]{extarticle}
\usepackage[a4paper,margin=1.5cm]{geometry}
\usepackage[T1]{fontenc}
\usepackage{amsmath,amssymb}
\usepackage{graphicx}
\usepackage[export]{adjustbox}
\usepackage{multicol}
\usepackage{enumitem}
\usepackage{booktabs,array,tabularx}
\setlength{\columnsep}{0.7cm}
\setlength{\parindent}{0pt}
\setlength{\parskip}{3pt}
\pagestyle{empty}
\makeatletter
\newcommand{\qfigdef}[4]{\@namedef{qf@file@#1}{#2}\@namedef{qf@w@#1}{#3}\@namedef{qf@wide@#1}{#4}}
\newcommand{\qf@check}[1]{\@ifundefined{qf@file@#1}{\PackageError{score2}{Unknown figure id #1}{}}{}}
\newcommand{\qf@img}[4][]{\includegraphics[width=#2\dimexpr\@nameuse{qf@w@#4}pt\relax,max width=#3,max height=0.45\textheight,#1]{\@nameuse{qf@file@#4}}}
\newif\ifqf@full
\newif\ifqf@boxed
\newsavebox{\qf@box}
\newcommand{\qf@open}{\global\qf@fulltrue\par\global\qf@boxedtrue
  \setbox\qf@box\vbox\bgroup\hsize\textwidth\linewidth\textwidth}
\newcommand{\qf@emit}[1]{\par\egroup\global\qf@boxedfalse
  \ifdim\dimexpr\ht\qf@box+\dp\qf@box\relax<#1\textheight\noindent\usebox{\qf@box}\par\else\unvbox\qf@box\fi}
\newcommand{\qfullsplit}{\ifqf@boxed\qf@emit{0.85}\fi}
\newcommand{\qf@close}{\ifqf@boxed\qf@emit{0.45}\fi\global\qf@fullfalse}
\newcommand{\qfullbegin}{\end{multicols}\qf@open}
\newcommand{\qfullend}{\qf@close\begin{multicols}{2}}
\newcommand{\qfullstart}{\qf@open}
\newcommand{\qfulldone}{\qf@close}
\newcommand{\qfig}[2][1]{\qf@check{#2}%
  \ifqf@full\qf@blk{#1}{#2}\else
  \ifnum\@nameuse{qf@wide@#2}=1 \expandafter\@firstoftwo\else\expandafter\@secondoftwo\fi
  {\end{multicols}\begin{center}\qf@img{#1}{\textwidth}{#2}\end{center}\begin{multicols}{2}}%
  {\qf@blk{#1}{#2}}\fi}
\newcommand{\qf@blk}[2]{\par\smallskip{\centering\qf@img{#1}{\linewidth}{#2}\par}\smallskip}
\newcommand{\qfigi}[2][1]{\qf@check{#2}\qf@img[valign=c]{#1}{\linewidth}{#2}}
\makeatother
\newcommand{\papersec}[1]{\par\bigskip{\centering\Large\bfseries PAPER-#1\par}\medskip}
\newcommand{\partsec}[2]{\par\medskip\noindent{\large\bfseries PART-#1: #2}\par\smallskip}
\newcommand{\sect}[1]{\par\smallskip\noindent\textbf{#1}\par}
\newcommand{\qn}[1]{\par\smallskip\noindent\textbf{#1.}~}
\newcommand{\paraq}[2]{\par\smallskip\noindent\textbf{Paragraph for Questions #1 and #2}\par}
\newenvironment{qoptions}{\begin{itemize}[leftmargin=1.8em,itemsep=1pt,topsep=2pt,label={}]}{\end{itemize}}
\newcommand{\opt}[2]{\item[(#1)] #2}
% figure registry (generated by score2_tools.py finalize)
\qfigdef{p1_phy_s1_q01_1}{SRG (Enthusiast) - Test 16 (12 Jan 2022)-img/p1_phy_s1_q01_1.png}{129.4}{0}
\qfigdef{p1_phy_s1_q02_1}{SRG (Enthusiast) - Test 16 (12 Jan 2022)-img/p1_phy_s1_q02_1.png}{154.5}{0}
\qfigdef{p1_phy_s1_q03_1}{SRG (Enthusiast) - Test 16 (12 Jan 2022)-img/p1_phy_s1_q03_1.png}{119.5}{0}
\qfigdef{p1_phy_s1_q05_1}{SRG (Enthusiast) - Test 16 (12 Jan 2022)-img/p1_phy_s1_q05_1.png}{87.9}{0}
\qfigdef{p1_phy_s1_q06_1}{SRG (Enthusiast) - Test 16 (12 Jan 2022)-img/p1_phy_s1_q06_1.png}{173.0}{0}
\qfigdef{p1_phy_s2_q07_1}{SRG (Enthusiast) - Test 16 (12 Jan 2022)-img/p1_phy_s2_q07_1.png}{73.5}{0}
\qfigdef{p1_phy_s2_q08_1}{SRG (Enthusiast) - Test 16 (12 Jan 2022)-img/p1_phy_s2_q08_1.png}{86.4}{0}
\qfigdef{p1_phy_s2_q09_2}{SRG (Enthusiast) - Test 16 (12 Jan 2022)-img/p1_phy_s2_q09_2.png}{72.5}{0}
\qfigdef{p1_phy_s2_q09_1}{SRG (Enthusiast) - Test 16 (12 Jan 2022)-img/p1_phy_s2_q09_1.png}{62.0}{0}
\qfigdef{p1_phy_s2_q10_1}{SRG (Enthusiast) - Test 16 (12 Jan 2022)-img/p1_phy_s2_q10_1.png}{140.7}{0}
\qfigdef{p1_phy_s2_q11_1}{SRG (Enthusiast) - Test 16 (12 Jan 2022)-img/p1_phy_s2_q11_1.png}{127.8}{0}
\qfigdef{p1_phy_s2_q11_2}{SRG (Enthusiast) - Test 16 (12 Jan 2022)-img/p1_phy_s2_q11_2.png}{95.1}{0}
\qfigdef{p1_phy_s2_q12_1}{SRG (Enthusiast) - Test 16 (12 Jan 2022)-img/p1_phy_s2_q12_1.png}{98.6}{0}
\qfigdef{p1_phy_s3_q14_1}{SRG (Enthusiast) - Test 16 (12 Jan 2022)-img/p1_phy_s3_q14_1.png}{96.2}{0}
\qfigdef{p1_phy_s3_q15_1}{SRG (Enthusiast) - Test 16 (12 Jan 2022)-img/p1_phy_s3_q15_1.png}{147.7}{0}
\qfigdef{p1_phy_s3_q18_1}{SRG (Enthusiast) - Test 16 (12 Jan 2022)-img/p1_phy_s3_q18_1.png}{149.9}{0}
\qfigdef{p1_chem_s1_q01_1}{SRG (Enthusiast) - Test 16 (12 Jan 2022)-img/p1_chem_s1_q01_1.png}{153.4}{0}
\qfigdef{p1_chem_s1_q01_3}{SRG (Enthusiast) - Test 16 (12 Jan 2022)-img/p1_chem_s1_q01_3.png}{225.8}{0}
\qfigdef{p1_chem_s1_q02_A}{SRG (Enthusiast) - Test 16 (12 Jan 2022)-img/p1_chem_s1_q02_A.png}{57.6}{0}
\qfigdef{p1_chem_s1_q02_B}{SRG (Enthusiast) - Test 16 (12 Jan 2022)-img/p1_chem_s1_q02_B.png}{60.2}{0}
\qfigdef{p1_chem_s1_q02_C}{SRG (Enthusiast) - Test 16 (12 Jan 2022)-img/p1_chem_s1_q02_C.png}{33.2}{0}
\qfigdef{p1_chem_s1_q02_D}{SRG (Enthusiast) - Test 16 (12 Jan 2022)-img/p1_chem_s1_q02_D.png}{33.2}{0}
\qfigdef{p1_chem_s1_q05_1}{SRG (Enthusiast) - Test 16 (12 Jan 2022)-img/p1_chem_s1_q05_1.png}{100.2}{0}
\qfigdef{p1_chem_s2_q09_1}{SRG (Enthusiast) - Test 16 (12 Jan 2022)-img/p1_chem_s2_q09_1.png}{184.8}{0}
\qfigdef{p1_chem_s2_q10_1}{SRG (Enthusiast) - Test 16 (12 Jan 2022)-img/p1_chem_s2_q10_1.png}{264.2}{1}
\qfigdef{p1_chem_s2_q10_2}{SRG (Enthusiast) - Test 16 (12 Jan 2022)-img/p1_chem_s2_q10_2.png}{68.7}{0}
\qfigdef{p1_chem_s2_q10_3}{SRG (Enthusiast) - Test 16 (12 Jan 2022)-img/p1_chem_s2_q10_3.png}{76.2}{0}
\qfigdef{p1_chem_s2_q10_5}{SRG (Enthusiast) - Test 16 (12 Jan 2022)-img/p1_chem_s2_q10_5.png}{76.2}{0}
\qfigdef{p1_chem_s2_q10_4}{SRG (Enthusiast) - Test 16 (12 Jan 2022)-img/p1_chem_s2_q10_4.png}{59.4}{0}
\qfigdef{p1_chem_s3_q17_1}{SRG (Enthusiast) - Test 16 (12 Jan 2022)-img/p1_chem_s3_q17_1.png}{370.9}{1}
\qfigdef{p2_phy_s1_q02_1}{SRG (Enthusiast) - Test 16 (12 Jan 2022)-img/p2_phy_s1_q02_1.png}{90.5}{0}
\qfigdef{p2_phy_s1_q03_1}{SRG (Enthusiast) - Test 16 (12 Jan 2022)-img/p2_phy_s1_q03_1.png}{87.5}{0}
\qfigdef{p2_phy_s1_q04_1}{SRG (Enthusiast) - Test 16 (12 Jan 2022)-img/p2_phy_s1_q04_1.png}{45.6}{0}
\qfigdef{p2_phy_s2_q05_1}{SRG (Enthusiast) - Test 16 (12 Jan 2022)-img/p2_phy_s2_q05_1.png}{70.9}{0}
\qfigdef{p2_phy_s2_q06_1}{SRG (Enthusiast) - Test 16 (12 Jan 2022)-img/p2_phy_s2_q06_1.png}{90.7}{0}
\qfigdef{p2_phy_s2_q11_1}{SRG (Enthusiast) - Test 16 (12 Jan 2022)-img/p2_phy_s2_q11_1.png}{173.3}{0}
\qfigdef{p2_chem_s1_q04_1}{SRG (Enthusiast) - Test 16 (12 Jan 2022)-img/p2_chem_s1_q04_1.png}{56.9}{0}
\qfigdef{p2_chem_s2_q06_1}{SRG (Enthusiast) - Test 16 (12 Jan 2022)-img/p2_chem_s2_q06_1.png}{69.6}{0}
\qfigdef{p2_chem_s2_q07_2}{SRG (Enthusiast) - Test 16 (12 Jan 2022)-img/p2_chem_s2_q07_2.png}{245.2}{0}
\qfigdef{p2_chem_s2_q08_1}{SRG (Enthusiast) - Test 16 (12 Jan 2022)-img/p2_chem_s2_q08_1.png}{171.3}{0}
\qfigdef{p2_chem_s2_q09_1}{SRG (Enthusiast) - Test 16 (12 Jan 2022)-img/p2_chem_s2_q09_1.png}{176.3}{0}
\qfigdef{p2_chem_s2_q09_A}{SRG (Enthusiast) - Test 16 (12 Jan 2022)-img/p2_chem_s2_q09_A.png}{113.3}{0}
\qfigdef{p2_chem_s2_q09_B}{SRG (Enthusiast) - Test 16 (12 Jan 2022)-img/p2_chem_s2_q09_B.png}{204.5}{0}
\qfigdef{p2_chem_s2_q09_C}{SRG (Enthusiast) - Test 16 (12 Jan 2022)-img/p2_chem_s2_q09_C.png}{162.4}{0}
\qfigdef{p2_chem_s2_q09_D}{SRG (Enthusiast) - Test 16 (12 Jan 2022)-img/p2_chem_s2_q09_D.png}{245.2}{0}
\qfigdef{p2_chem_s3_q15_1}{SRG (Enthusiast) - Test 16 (12 Jan 2022)-img/p2_chem_s3_q15_1.png}{326.0}{1}
\qfigdef{p2_chem_s3_q16_1}{SRG (Enthusiast) - Test 16 (12 Jan 2022)-img/p2_chem_s3_q16_1.png}{320.3}{1}
\qfigdef{p2_math_s2_q10_1}{SRG (Enthusiast) - Test 16 (12 Jan 2022)-img/p2_math_s2_q10_1.png}{95.8}{0}
\qfigdef{p2_math_s3_q17_1}{SRG (Enthusiast) - Test 16 (12 Jan 2022)-img/p2_math_s3_q17_1.png}{106.7}{0}

\begin{document}
\begin{center}{\Large\bfseries SRG (Enthusiast) --- Test 16 (12 Jan 2022)}\end{center}
\begin{multicols}{2}
\papersec{1}
\partsec{1}{PHYSICS}
\sect{Section-I(i)}
\qn{1} ABCD is a square of side $4R$ and EFGH is a square loop of side $2R$. The loop carries a current $I$. The magnetic field at the centre O of the loop
\qfig{p1_phy_s1_q01_1}
\begin{qoptions}
\opt{A}{$\dfrac{\mu_0 I}{2\pi R}\ \odot$}
\opt{B}{$\dfrac{\mu_0 I}{2\pi R}\ \otimes$}
\opt{C}{$\dfrac{\mu_0 I}{\sqrt2\,2\pi R}\ \otimes$}
\opt{D}{$\dfrac{\mu_0 I}{\sqrt2\,2\pi R}\ \odot$}
\end{qoptions}
\qn{2} A thin lens with radius of curvature 30 cm and 60 cm is placed with water on right side and air on left side, then the focal length of the lens for the parallel incident rays shown in the figure, is
\qfig{p1_phy_s1_q02_1}
\begin{qoptions}
\opt{A}{72 cm}
\opt{B}{144 cm}
\opt{C}{36 cm}
\opt{D}{None of these}
\end{qoptions}
\qn{3} Young's double slit experiment is submerged in water as shown in the figure. If an air bubble is formed on slit $S_1$, then the fringe pattern
\qfig{p1_phy_s1_q03_1}
\begin{qoptions}
\opt{A}{Will shift towards A}
\opt{B}{Will shift towards B}
\opt{C}{Will not shift}
\opt{D}{Can't say}
\end{qoptions}
\qn{4} Magnetic field exist in the space and given as $\vec B=-\dfrac{B_0}{l^2}x^2\hat k$ where $B_0$ and $l$ are positive constants. A particle having positive charge `$q$' and mass `$m$' is projected with speed `$v_0$' along positive x-axis from the origin. What is the maximum distance of the charged particle from the y-axis before it turns back due to the magnetic field. (Ignore any interaction other than magnetic field)
\begin{qoptions}
\opt{A}{$\left(\dfrac{ml^2v_0}{3qB_0}\right)^{1/3}$}
\opt{B}{$\left(\dfrac{3ml^2v_0}{2qB_0}\right)^{1/3}$}
\opt{C}{$\left(\dfrac{3ml^2v_0}{qB_0}\right)^{1/3}$}
\opt{D}{$\left(\dfrac{ml^2v_0}{qB_0}\right)^{1/3}$}
\end{qoptions}
\qn{5} A non-conducting hollow cone has charge density $\sigma$. A part ABP is cut and removed from the cone. The potential due to the remaining portion of the cone at point `P' is
\qfig{p1_phy_s1_q05_1}
\begin{qoptions}
\opt{A}{$\dfrac{5}{6}\dfrac{\sigma R}{\varepsilon_0}$}
\opt{B}{$\dfrac{5}{24}\dfrac{\sigma R}{\varepsilon_0}$}
\opt{C}{$\dfrac{5}{3}\dfrac{\sigma R}{\varepsilon_0}$}
\opt{D}{$\dfrac{5}{12}\dfrac{\sigma R}{\varepsilon_0}$}
\end{qoptions}
\qn{6} A point charge `$q$' is placed at a distance `$l$' (on the y-axis) vertically above the surface which lies in the xz plane as shown in the figure. The flux of electric field passing through the plate is
\qfig{p1_phy_s1_q06_1}
\begin{qoptions}
\opt{A}{$\dfrac{q}{6\varepsilon_0}$}
\opt{B}{$\dfrac{q}{18\varepsilon_0}$}
\opt{C}{$\dfrac{q}{24\varepsilon_0}$}
\opt{D}{$\dfrac{q}{48\varepsilon_0}$}
\end{qoptions}
\sect{Section-I(ii)}
\qn{7} In the circuit shown in the figure $C_1=C_2=C$ and capacitor $C_1$ is having initial charge $Q$. The switch is closed at $t=0$. Which of the following options are correct
\qfig{p1_phy_s2_q07_1}
\begin{qoptions}
\opt{A}{The current at time $t$ through the circuit will be $\left(\dfrac{Q}{RC}\right)e^{-2t/RC}$}
\opt{B}{The charge on capacitor $C_1$ at time $t$ will be $\dfrac{Q}{2}\left(1+e^{-2t/RC}\right)$}
\opt{C}{The charge on the capacitor $C_2$ at time $t$ will be $\dfrac{Q}{2}\left(1-e^{-2t/RC}\right)$}
\opt{D}{Total heat loss in the process of charge transfer will be $\dfrac{Q^2}{8C}$}
\end{qoptions}
\qn{8} Four identical mirrors are made to stand vertically to form a square arrangement as shown in the top view. A ray starts from the midpoint M of mirror AD and after two reflections reached corner D. Then, angle $\theta$ must be
\qfig{p1_phy_s2_q08_1}
\begin{qoptions}
\opt{A}{$\tan^{-1}(0.75)$}
\opt{B}{$\cot^{-1}(0.75)$}
\opt{C}{$\sin^{-1}(0.75)$}
\opt{D}{$\cos^{-1}(0.75)$}
\end{qoptions}
\qn{9} In YDSE, intensity of light wave passing through slits $S_1$ and $S_2$ are $I_1$ and $I_2$ respectively as shown in figure (1). The pattern of intensity distribution on the screen is given in figure (2) where $x$ is distance measured from the point O. In this case which of the following statement(s) is / are true\\
(Given $d\ll D$)
\qfig{p1_phy_s2_q09_2}
{\centering Figure 1\par}
\qfig{p1_phy_s2_q09_1}
{\centering Figure 2\par}
\begin{qoptions}
\opt{A}{The intensity of light passing through $S_1$ and $S_2$ is the same}
\opt{B}{The phase difference between the light waves passing through $S_1$ and $S_2$ to a fixed point on the screen is independent of time}
\opt{C}{The phase difference between the light waves passing through $S_1$ and $S_2$ to a fixed point on the screen varies with time}
\opt{D}{Light waves passing through $S_1$ and $S_2$ have the same wave length}
\end{qoptions}
\qn{10} In the circuit shown in the figure $I_1,I_2,I_3,I_4$ denote current in the branches. Choose the \textbf{INCORRECT} option
\qfig{p1_phy_s2_q10_1}
\begin{qoptions}
\opt{A}{$I_1>I_2>I_3>I_4$}
\opt{B}{$I_1=I_2>I_3>I_4$}
\opt{C}{$I_1<I_3<I_2<I_4$}
\opt{D}{$I_2<I_1<I_3<I_4$}
\end{qoptions}
\qn{11} In the figure Box-1 contains one capacitor and one resistor in series
\qfig{p1_phy_s2_q11_1}
Box - 2 contains one inductor and one resistor in series\\
$R_1$ = resistance in Box 1\\
$R_2$ = resistance in Box 2\\
Power factor of Box 1 $=\dfrac{1}{\sqrt5}$\\
Power factor of Box 2 $=\dfrac{1}{\sqrt5}$\\
Power factor of (Box 1 + Box 2) $=\dfrac{1}{\sqrt2}$
\qfig{p1_phy_s2_q11_2}
Power factor of Box 3 $=\dfrac{1}{\sqrt2}$\\
Resistance of Box 3 $=R_1+R_2$\\
Current $I=I_0\sin(\omega t+\phi)$\\
On the basis of above mentioned information pick the correct option(s) from below :
\begin{qoptions}
\opt{A}{$\dfrac{R_1}{R_2}=3$}
\opt{B}{$Z_3=4\sqrt2R_2$ (Impedance of Box 3)}
\opt{C}{$I=\dfrac{V_0}{4R_2}\sin(\omega t)$}
\opt{D}{$\phi=(2n+1)\pi$, where $n$ may be some integer}
\end{qoptions}
\qn{12} In the figure ABCDEFA is a closed conducting wire loop. The part AFE is a portion of an ellipse $\dfrac{X^2}{4a^2}+\dfrac{Y^2}{a^2}=1$ and BCD is a portion of the ellipse $\dfrac{X^2}{9a^2}+\dfrac{Y^2}{4a^2}=1$. A uniform magnetic field $\vec B=B_0\hat k$ exists in this region. If the frame is rotated about the y-axis with a constant angular velocity `$\omega$'. The maximum emf induced in the loop is ,
\qfig{p1_phy_s2_q12_1}
\begin{qoptions}
\opt{A}{$\dfrac{\pi\omega B_0a^2}{2}$}
\opt{B}{$\pi\omega B_0a^2$}
\opt{C}{$2\pi\omega B_0a^2$}
\opt{D}{$1.5\pi\omega B_0a^2$}
\end{qoptions}
\sect{Section-II}
\qn{13} Wavelength of the $k_\alpha$ characteristics X-ray of iron and potassium are $1.931\times10^{-8}$ cm and $3.737\times10^{-8}$ cm respectively. The atomic number (approximately) of the element for which characteristic $k_\alpha$ wavelength is $2.289\times10^{-8}$ cm is $k$, the value of $10k$ is\\
(atomic number of iron is 26 and atomic number is potassium is 19)
\qn{14} A charged particle A hangs with an inextensible string of length $l$. Another charged particle B is fixed at a distance $l=10$ m below the point of suspension O as shown in the figure. In the equilibrium the particle A is at a height $h=3$ m above the level of particle B. If the mass of the particle A is $m=2$ kg and acceleration due to gravity is $g=10\,\mathrm{m/s^2}$, then the electrical potential energy of interaction of the particle in the equilibrium will be
\qfig{p1_phy_s3_q14_1}
\qn{15} If final image formed after two refractions through the lens and one reflection from the mirror forming at same point o then find maximum value of `d'. (if $\mu$ is refractive index of material of lens is 1.5)
\qfig{p1_phy_s3_q15_1}
\qn{16} A ray of light falling on a glass sphere of refractive index $\mu=\sqrt3$ such that the directions of the incident ray and emergent ray when produced meet the surface at the same point. The value of angle of incidence is $\theta$. Calculate the value of $2\theta$ in degrees.
\qn{17} A charged particle having charge $q$ and mass $m$ is projected into a region of uniform electric field of strength $\vec E_0$, with velocity $\vec V_0$ perpendicular to $\vec E_0$. Throughout the motion apart from electric force particle also experiences a dissipative force of constant magnitude $qE_0$ and directed opposite to its velocity. If $|\vec V_0|=6$ m/s, then its speed when it has turned through an angle of $90^\circ$, is equal to
\qn{18} Consider a plane inclined at an angle $45^\circ$ with the horizontal has two slits ($S_1$ and $S_2$) separated by a distance $d=\sqrt2$ mm. The screen is placed at a distance of $D=10$ m. A parallel monochromatic light beam of wavelength 5000 \AA\ is incident on the slits as shown. If the fringe width of interference pattern on the screen is $k\times10^{-3}$ meter, then find the value of $k$.
\qfig{p1_phy_s3_q18_1}
\partsec{2}{CHEMISTRY}
\sect{Section-I(i)}
\qn{1}
\qfig{p1_chem_s1_q01_1}
To carry out above conversion, arrange the following reagents in correct order\\
$\mathrm{O_3/Zn,\Delta}$\quad $\mathrm{dil.OH^{\ominus}/\Delta}$\quad NaOCl\quad $\mathrm{H^+}$\\
(1)\qquad\qquad (2)\qquad\qquad (3)\qquad (4)
\begin{qoptions}
\opt{A}{$1\to3\to2\to4$}
\opt{B}{$1\to2\to4\to3$}
\opt{C}{$1\to3\to4\to2$}
\opt{D}{$1\to2\to3\to4$}
\end{qoptions}
\qn{2}
\qfig{p1_chem_s1_q01_3}
Product (A) of the above reaction is :
\begin{qoptions}
\opt{A}{\qfigi{p1_chem_s1_q02_A}}
\opt{B}{\qfigi{p1_chem_s1_q02_B}}
\opt{C}{\qfigi{p1_chem_s1_q02_C}}
\opt{D}{\qfigi{p1_chem_s1_q02_D}}
\end{qoptions}
\qn{3} In which of the following complexes bond order of C--O bond is minimum ?
\begin{qoptions}
\opt{A}{$\left[\mathrm{Mn(CO)_6}\right]^{+}$}
\opt{B}{$\left[\mathrm{Cr(CO)_6}\right]$}
\opt{C}{$\left[\mathrm{Ti(CO)_6}\right]^{2-}$}
\opt{D}{$\left[\mathrm{V(CO)_6}\right]^{-}$}
\end{qoptions}
\qn{4} Identify M and N in the following reaction.\\
Copper glance $\xrightarrow[\text{in presence of air}]{\text{controlled heating}}\mathrm{M+SO_2}$\\
$\mathrm{M}\xrightarrow{N}\mathrm{Cu+SO_2}\uparrow$
\begin{qoptions}
\opt{A}{$\mathrm{M=Cu_2O}$ ; N = Self reduction}
\opt{B}{$\mathrm{M=Cu_2O+Cu_2S}$ ; N = Only heating}
\opt{C}{$\mathrm{M=Cu_2O}$ ; N = Carbon reduction}
\opt{D}{$\mathrm{M=Cu_2O}$ ; N = Electrolytic reduction}
\end{qoptions}
\qn{5} What will be the order of reaction and rate constant for a chemical change having $\log\left[t_{50\%}\right]$ vs $\log$ [concentration (a)] curves as ?
\qfig{p1_chem_s1_q05_1}
\begin{qoptions}
\opt{A}{0, 1}
\opt{B}{1, 1}
\opt{C}{2, 2}
\opt{D}{3, 1}
\end{qoptions}
\qn{6} The cell reaction for the given cell is spontaneous if :\\
$\mathrm{Pt\ ;\ Cl_2}\ \big|\,\mathrm{Cl^-}\,\big\|\ \mathrm{Cl^-}\,\big|\ \mathrm{Cl_2\ ,\ Pt}$\\
\hspace*{1.6em}$P_1$\hspace{1.3em}1M\hspace{1.2em}1M\hspace{1em}$P_2$
\begin{qoptions}
\opt{A}{$P_1>P_2$}
\opt{B}{$P_1<P_2$}
\opt{C}{$P_1=P_2$}
\opt{D}{$P_2=1$ atm}
\end{qoptions}
\sect{Section-I(ii)}
\qn{7} Which of the following is / are correct for lyophilic aquasols ?
\begin{qoptions}
\opt{A}{Its surface tension is lower than that of $\mathrm{H_2O}$}
\opt{B}{Its viscosity is higher than that of water}
\opt{C}{Colloidal silver}
\opt{D}{Colloidal antimony}
\end{qoptions}
\qn{8} Nitrogen dioxide can be obtained by heating
\begin{qoptions}
\opt{A}{$\mathrm{KNO_3}$}
\opt{B}{$\mathrm{Pb(NO_3)_2}$}
\opt{C}{$\mathrm{LiNO_3}$}
\opt{D}{$\mathrm{AgNO_3}$}
\end{qoptions}
\qn{9} Which of the following statements are true about the major product (X) in the reaction ?
\qfig{p1_chem_s2_q09_1}
\begin{qoptions}
\opt{A}{(X) is a cyclic amide}
\opt{B}{(X) has an asymmetric centre}
\opt{C}{(X) is a primary amine}
\opt{D}{(X) can also be obtained by treatment of $\mathrm{NH_3}$ with $\mathrm{C_6H_5CH(CH_3)Br}$}
\end{qoptions}
\qfullbegin
\qn{10} In a reaction sequence
\qfig{p1_chem_s2_q10_1}
\qfullsplit
\begin{qoptions}
\opt{A}{A is \qfigi{p1_chem_s2_q10_2}}
\opt{B}{B is \qfigi{p1_chem_s2_q10_3}}
\opt{C}{A is \qfigi{p1_chem_s2_q10_5}}
\opt{D}{C is \qfigi{p1_chem_s2_q10_4}}
\end{qoptions}
\qfullend
\qn{11} Which of the following is / are correct ?
\begin{qoptions}
\opt{A}{On strong heating $\mathrm{FeCl_3}$ gives $\mathrm{FeCl_2}$ and $\mathrm{Cl_2}$}
\opt{B}{Dilute $\mathrm{CuCl_2}$ is blue due to formation of $\left[\mathrm{CuCl_4}\right]^{2-}$}
\opt{C}{$\mathrm{PH_3}$ can be obtained by heating red P with NaOH}
\opt{D}{$\mathrm{PH_3,H_3PO_3}$ and $\mathrm{H_3PO_2}$, all are reducing agents}
\end{qoptions}
\qn{12} $2X(g)+Y(g)+3Z(g)\to\mathrm{Products}$\\
The rate equation of the above reaction is given by :\\
Rate $=K[X][Y]^0[Z]^2$\\
Choose correct statements :
\begin{qoptions}
\opt{A}{The rate of reaction becomes doubled by doubling the conc.of (X)}
\opt{B}{Rate of reaction decreases by reducing the conc of (Y) to half of the original value}
\opt{C}{The half life of (Z) increases by increasing its conc}
\opt{D}{If 75\% of (X) undergoes reaction in 20 sec, 50\% of (X) will react in 10 sec}
\end{qoptions}
\sect{Section-II}
\qn{13} Freezing point of 0.2 M KCN (100\% ionisation) solution is $-0.7^\circ$C. On adding 0.1 mole of $\mathrm{Hg(CN)_2}$ to one litre of the 0.2 M KCN solution, the freezing point of the solution becomes $-0.53^\circ$C due to the reaction\\
$\mathrm{Hg(CN)_2+mCN^-\to Hg(CN)_{m+2}^{m-}}$. What is the value of $m$, assuming molality $=$ molarity ?
\qn{14} If $E^0_{\mathrm{Fe^{2+}/Fe}}=5\,$V and $E^0_{\mathrm{Fe^{3+}/Fe^{2+}}}=2\,$V then $E^0_{\mathrm{Fe^{3+}/Fe}}$ will be
\qn{15} A cell made up of two hydrogen electrodes. The positive electrode is in contact with a solution of $10^{-4}\,M$ hydrogen ions. The emf of the cell is 0.118 V at 25$^\circ$C. The concentration of hydrogen ions at negative electrode is $10^{-x}\,M$. Then $x=$ ?
\qn{16} When $\mathrm{Mg_2C_3}$ is reacted with water, a compound X is formed which on heating with Na followed by treatment with $\mathrm{C_2H_5Br}$ gives another compound Y. Calculate the number of sigma bond in Y.
\qfullbegin
\qn{17} How many of the following acids easily undergo decarboxylation on heating ?
\qfig{p1_chem_s3_q17_1}
\qfullend
\qn{18} No.of geometrical isomers for $\left[\mathrm{M(gly)Cl_2Br_2}\right]^{+}$ is :
\partsec{3}{MATHEMATICS}
\sect{Section-I(i)}
\qn{1} If $f:[0,1]\to[0,\infty)$ is a differentiable function with decreasing first derivative, such that $f(0)=0$ \& $f'(x)>0$ then
\begin{qoptions}
\opt{A}{$\displaystyle\int_0^1\dfrac{dx}{\left(f(x)\right)^2+1}>\dfrac{f(1)}{f'(1)}$}
\opt{B}{$\displaystyle\int_0^1\dfrac{dx}{\left(f(x)\right)^2+1}<\dfrac{f(1)}{f'(1)}$}
\opt{C}{$\displaystyle\int_0^1\dfrac{dx}{\left(f(x)\right)^2+1}>\dfrac{\tan^{-1}\left(f(1)\right)}{f'(1)}$}
\opt{D}{$\displaystyle\int_0^1\dfrac{dx}{\left(f(x)\right)^2+1}=\dfrac{f(1)}{f'(1)}$}
\end{qoptions}
\qn{2} Let $a\in[0,4]$. The maximum area bounded by the curves $y=1-|x-1|$ \& $y=|2x-a|$ is
\begin{qoptions}
\opt{A}{1/3}
\opt{B}{1/4}
\opt{C}{1/2}
\opt{D}{None of these}
\end{qoptions}
\qn{3} If $f:R\to R$ is a continuous function satisfying $f(0)=1$ and $f(2x)-f(x)=x\ \forall x\in R$ and $\displaystyle\lim_{n\to\infty}\left(f(x)-f\!\left(\dfrac{x}{2^n}\right)\right)=p(x)$ then $p(x)$ is
\begin{qoptions}
\opt{A}{Constant function}
\opt{B}{A linear polynomial in $x$}
\opt{C}{A quadratic polynomial in $x$}
\opt{D}{A cubic polynomial in $x$}
\end{qoptions}
\qn{4} If $f(x)$ is continuous function such that $f(x)>0\ \forall x\ge0$ and $\left(f(x)\right)^{2014}=1+\displaystyle\int_0^x f(t)\,dt$ then the value of $\left(f(2014)\right)^{2013}$ is equal to
\begin{qoptions}
\opt{A}{2013}
\opt{B}{2014}
\opt{C}{2015}
\opt{D}{2016}
\end{qoptions}
\qn{5} Let $f:[0,\infty)\to R$ be a continuous and strictly increasing function such that $f^4(x)=\displaystyle\int_0^x t^2f^3(t)\,dt\ \forall x>0$. The area enclosed by $y=f(x)$, the x-axis and the ordinates at $x=-1$ and $x=4$ is
\begin{qoptions}
\opt{A}{$\dfrac{17}{48}$}
\opt{B}{$\dfrac{65}{48}$}
\opt{C}{$\dfrac{257}{48}$}
\opt{D}{$\dfrac{235}{48}$}
\end{qoptions}
\qn{6} For all $x\in\left[-\dfrac72,\dfrac72\right]$, let $f(x)=\begin{cases}\big|[x]\big|, & 0\le\{x\}<\dfrac12\\[4pt] |x|, & \dfrac12\le\{x\}<1\end{cases}$\\
Note : $[x]$ denotes greatest integer less than or equal to $x$ and $\{x\}$ denotes fractional part of $x$\\
If $L$ is the number of points of discontinuity and $M$ is the number of points of non-differentiability of the function $f(x)$, then $(L+M)$ is equal to
\begin{qoptions}
\opt{A}{18}
\opt{B}{20}
\opt{C}{21}
\opt{D}{22}
\end{qoptions}
\sect{Section-I(ii)}
\qn{7} Let $f,\ g$ \& $h$ are functions differentiable on some open interval around origin $(0,0)$ and satisfy the equations $f'=2f^2gh+\dfrac{1}{gh}$, $f(0)=1$, $g'=fg^2h+\dfrac{4}{fh}$, $g(0)=1$, $h'=3fgh^2+\dfrac{1}{fg}$, $h(0)=1$.\\
The function $f$ is given by
\begin{qoptions}
\opt{A}{$f(x)=2^{-\frac{1}{12}}\left(\dfrac{\sin\left(6x+\frac{\pi}{4}\right)}{\cos^2\left(6x+\frac{\pi}{4}\right)}\right)^{1/6}$}
\opt{B}{$f(x)=2^{-\frac{1}{6}}\left(\dfrac{\sin\left(6x+\frac{\pi}{4}\right)}{\cos^2\left(6x+\frac{\pi}{4}\right)}\right)^{1/12}$}
\opt{C}{$f(x)=(\sec12x)^{1/12}(\sec12x+\tan12x)^{1/4}$}
\opt{D}{$f(x)=(\sec12x)^{1/14}(\sec12x+\tan12x)^{1/12}$}
\end{qoptions}
\qn{8} For any real number $\lambda\ge1$. Let $f(\lambda)$ denote the real solution of the equation $x(1+\ln x)=\lambda$ then
\begin{qoptions}
\opt{A}{$f(\lambda)$ is strictly increasing function in its domain}
\opt{B}{$\displaystyle\lim_{\lambda\to\infty}f(\lambda)\to\infty$}
\opt{C}{If $h(x)=x(1+\ln x)$, $h:[1,\infty)\to[1,\infty)$ then $f(\lambda)$ is inverse function of $h(x)$}
\opt{D}{$\displaystyle\lim_{\lambda\to\infty}\dfrac{f(\lambda)\ln\lambda}{\lambda}=1$}
\end{qoptions}
\qn{9} A point $Q$ moves on the positive y-axis and a point $P$ follows the point $Q$ such that the motion of $P$ is always towards $Q$ and distance between $P$ and $Q$ is always equal to some constant $K>0$. If initially $Q$ is at origin and $P$ is at $(K,0)$ then
\begin{qoptions}
\opt{A}{The differential equation representing the curve traced by $P$ is $\dfrac{dy}{dx}=\dfrac{\sqrt{K^2-x^2}}{-x}$}
\opt{B}{y-axis is asymptote to the curve traced by $P$}
\opt{C}{The equation of curve traced by $P$ is $y=-K\ln\left(\dfrac{K+\sqrt{K^2-x^2}}{x}\right)-\sqrt{K^2-x^2}$}
\opt{D}{The curve traced by $P$ will have x-axis as asymptote}
\end{qoptions}
\qn{10} The function $f(x)=x^{\frac13}(x-1)$
\begin{qoptions}
\opt{A}{Has vertical tangent at $x=0$}
\opt{B}{Area enclosed by curve with x-axis from $x=-1$ to $x=1$ is $\dfrac32$}
\opt{C}{Is concave up in $(0,\infty)$}
\opt{D}{Is decreasing for $x\in\left(-\infty,\dfrac14\right)$}
\end{qoptions}
\qn{11} Let $\vec r$ be a position vector of a variable point in cartesian xy-plane such that $\vec r\cdot\left(-8\hat i+10\hat j-\vec r\right)=40$ and $M=\max\left|\vec r+2\hat i-3\hat j\right|^2,\ m=\min\left|\vec r+2\hat i-3\hat j\right|^2$, then $(M+m)$ is less than or equal to
\begin{qoptions}
\opt{A}{9}
\opt{B}{18}
\opt{C}{27}
\opt{D}{28}
\end{qoptions}
\qn{12} A computer producing factory has only two plants $T_1$ and $T_2$. Plant $T_1$ produces 20\% and plant $T_2$ produces 80\% of the total computers produced. 7\% of computers produced in the factory turn out to be defective. It is known that\\
$P$(computer turns out to be defective given that it is produced in plant $T_1$)\\
\hspace*{1.5em}$=10\,P$(computer turns out to be defective given that it is produced in plant $T_2$),\\
where $P(E)$ denotes the probability of an event $E$. A computer produced in the factory is randomly selected and it does not turn out to be defective. Then the probability that it is produced in plant $T_2$ is
\begin{qoptions}
\opt{A}{$\dfrac{36}{73}$}
\opt{B}{$\dfrac{47}{79}$}
\opt{C}{$\dfrac{78}{93}$}
\opt{D}{$\dfrac{75}{83}$}
\end{qoptions}
\sect{Section-II}
\qn{13} If $\displaystyle\lim_{r\to\infty}\dfrac{r^C\displaystyle\int_0^{\pi/2}x^r\sin x\,dx}{\displaystyle\int_0^{\pi/2}x^r\cos x\,dx}=L$, then value of $\pi L-C$ is equal to \dotfill $(C\in R,\ L>0)$
\qn{14} $n$ biased coins, with $m^{\text{th}}$ coin having probability of throwing head equal to $\dfrac{1}{(2m+1)}$ ($m=1,2,\dots\dots,n$), are tossed once. The probability of getting an odd numbers of heads, if results for each coin are independent, is $\dfrac{\alpha n}{\beta n+\gamma}$ then $\alpha+\beta+\gamma$ is equal to (given $\alpha,\beta\ \&\ \gamma$ are co-prime natural numbers)
\qn{15} Let $\vec A=\hat i-2\hat j+3\hat k,\ \vec B=2\hat i+\hat j-\hat k,\ \vec C=\hat j+\hat k$. If the vector $\vec B\times\vec C$ can be expressed as a linear combination $\vec B\times\vec C=x\vec A+y\vec B+z\vec C$ where $x,y,z$ are scalar, then find the value of $x+y+z$ \dots\dots.
\qn{16} Find the sum of an infinite geometric series whose first term is $\displaystyle\lim_{x\to0}\sum_{k=1}^{2013}\dfrac{\left\{\dfrac{x}{\tan x}+2k\right\}}{2013}$\\
and whose common ratio is the value of $\displaystyle\lim_{x\to0}\dfrac{e^{\tan^3x}-e^{x^3}}{2\ln\left(1+x^3\sin^2x\right)}$\\
{[}Note : where $\{k\}$ denotes fractional part of $k$]
\qn{17} If $\vec V_1=\hat i-2\hat j+3\hat k$ and $\vec V_2=a\hat i+b\hat j+c\hat k$ where $a,b,c\in\{-1,0,1,2,3\}$ then find the number of non-zero vectors $\vec V_2$ which are normal to $\vec V_1$
\qn{18} Let $A=\begin{pmatrix}2x\\5\\-3x^2\end{pmatrix}$, $B=\begin{pmatrix}2a&-b&5c\end{pmatrix}$, $C=\begin{pmatrix}5x&-1&4x^2\end{pmatrix}$ and $D=\begin{pmatrix}2\\1\\4\end{pmatrix}$ be four given matrices, where $a,b,c$ and $x\in R$.\\
Given that $t_r(AB)=t_r(DC)\ \forall x\in R$ where $t_r(A)$ denotes trace of matrix $A$, then if $\vec V_1=a\hat i+b\hat j+c\hat k,\ \vec V_2=\dfrac{2\hat i}{5}+10\hat j-30\hat k$ and $\vec V_3=\hat k$ then find the value of $\left[\left(\vec V_1\times\vec V_2\right)\cdot\vec V_3\right]^2$\\
{[}Note : $[k]$ denotes greatest integer less than or equal to $k$]
\papersec{2}
\partsec{1}{PHYSICS}
\sect{Section-I(i)}
\qn{1} You are holding a bottle of bubbled water inside a car moving forward. When the driver applies the brakes:
\begin{qoptions}
\opt{A}{Bubbles in the middle of the liquid will start to move forward with respect to the bottle.}
\opt{B}{Bubbles will start to move backward with respect to the bottle.}
\opt{C}{Bubbles will stay at the same horizontal location in the water}
\opt{D}{Depending on the speed of the car, bubbles might move forward or backward}
\end{qoptions}
\qn{2} A horizontal force $F$ is applied at the top of an equilateral triangular block having mass $m$. The minimum coefficient of friction required to topple the block before translation will be
\qfig{p2_phy_s1_q02_1}
\begin{qoptions}
\opt{A}{$\dfrac{2}{\sqrt3}$}
\opt{B}{$\dfrac12$}
\opt{C}{$\dfrac{1}{\sqrt3}$}
\opt{D}{None of these}
\end{qoptions}
\qn{3} Four rods of equal lengths AB, BC, CD and DA have mass $m$, $2m$, $3m$ and $4m$ respectively are place in x-y plane (see figure) given $AB=BC=CD=DA=4$ meters. The centre of mass of the system will satisfy ($x$ and $y$ are in meters)
\qfig{p2_phy_s1_q03_1}
\begin{qoptions}
\opt{A}{$x>2,\ y>2$}
\opt{B}{$x>2,\ y<2$}
\opt{C}{$x<2,\ y<2$}
\opt{D}{$x<2,\ y>2$}
\end{qoptions}
\qn{4} A hollow non conducting sphere of radius $R$ has a positive charge $q$ uniformly distributed on its surface. The sphere starts rotating with a constant angular velocity $\omega$ about an axis passing through center of sphere, as shown in the figure. Then the net magnetic field at center of the sphere is
\qfig{p2_phy_s1_q04_1}
\begin{qoptions}
\opt{A}{$\dfrac{\mu_0q\omega}{2\pi R}$}
\opt{B}{$\dfrac{\mu_0q\omega}{3\pi R}$}
\opt{C}{$\dfrac{\mu_0q\omega}{6\pi R}$}
\opt{D}{$\dfrac{\mu_0q\omega}{12\pi R}$}
\end{qoptions}
\sect{Section-I(ii)}
\qn{5} In the circuit shown $R_1=R_2=10\,\Omega$ and resistance per unit length of wire PQ $=1\,\Omega$/cm and length PQ $=10$ cm. If $R_2$ is made $20\,\Omega$. Then to get zero deflection in galvanometer. S is midpoint of wire PQ
\qfig{p2_phy_s2_q05_1}
\begin{qoptions}
\opt{A}{The jockey at P can be moved towards right 2 cm}
\opt{B}{The jockey at Q can be moved towards left 2cm}
\opt{C}{The jockey at S can be moved towards left a distance 5/3 m}
\opt{D}{The jockey at all position fixed and $R_1$ should be made $20\,\Omega$}
\end{qoptions}
\qn{6} In a series LCR circuit, different physical quantities vary with frequency $\omega$. Which of the following curves represent correct frequency variation of the corresponding quantity?
\qfig{p2_phy_s2_q06_1}
\begin{qoptions}
\opt{A}{Curve I for $R$}
\opt{B}{Curve II for current $I$}
\opt{C}{Curve III for $X_L$}
\opt{D}{Curve IV for $X_c$}
\end{qoptions}
\qn{7} An observer `O' at an accelerated trolley observes a projectile Q of mass 5 kg moving under uniform gravitational field of earth acting towards negative z-axis. He observes the momentum $\vec p$ of the projectile Q as $\vec P=15t\hat i-10t^2\hat j-50t\hat k$ (kg m/s), where $t$ denotes time in second. Find the acceleration of the trolley at $t=2$ sec. (Take $g=10$ m/s$^2$)
\begin{qoptions}
\opt{A}{$6\hat i-8\hat j-20\hat k\,\mathrm{m/s^2}$}
\opt{B}{$3\hat i-8\hat j\,\mathrm{m/s^2}$}
\opt{C}{$-3\hat i-8\hat j\,\mathrm{m/s^2}$}
\opt{D}{$-3\hat i+8\hat j\,\mathrm{m/s^2}$}
\end{qoptions}
\qn{8} Pick the correct option
\begin{qoptions}
\opt{A}{Elastic deformation energy is basically electrostatic potential energy.}
\opt{B}{In general the angle of contact of a solid liquid pair decreases with rise in temperature.}
\opt{C}{Siphon action can not be explained by capillary rise phenomenon.}
\opt{D}{An ideal fluid must be transparent.}
\end{qoptions}
\qn{9} X-rays from a tube with a target A of atomic number $Z$ shows strong K line for target A and weak K lines for impurities. The wavelength of $K_\alpha$ lines is $\lambda_z$ for target A and $\lambda_1$ and $\lambda_2$ for two impurities. If $\dfrac{\lambda_z}{\lambda_1}=4$ and $\dfrac{\lambda_z}{\lambda_2}=\dfrac14$ and taking screening constant of $K_\alpha$ lines to be unity, select the correct statement(s)
\begin{qoptions}
\opt{A}{The atomic number of first impurity is $2Z-1$}
\opt{B}{The atomic number of first impurity is $2Z+1$}
\opt{C}{The atomic number of second impurity is $\dfrac{(Z+1)}{2}$}
\opt{D}{The atomic number of second impurity is $\dfrac Z2+1$}
\end{qoptions}
\qn{10} An equi-convex lens ($\mu=1.5$) is combined with an equi concave lens ($\mu=1.3$). The radii of curvature of all surfaces is the same. (Assume thin lenses)
\begin{qoptions}
\opt{A}{The combination behaves like a converging lens in medium of refractive index 1.4 and diverging lens in a medium of refractive index 1.6.}
\opt{B}{The combination behaves like a converging lens in a medium of refractive index 1.2 as well as 1.4.}
\opt{C}{The combination behaves like a diverging lens in medium of refractive index 1.2 as well as 1.6.}
\opt{D}{The combination behaves like a diverging lens in medium of refractive index 1.4 and converging lens in a medium of refractive index 1.6.}
\end{qoptions}
\qn{11} A block of mass 1 gm carrying a charge $1\,\mu$C is connected to a spring of stiffness $\left(\pi^2\times10^{-3}\right)$N/m whose other end is fixed to a wall. The block rests on a smooth ground as shown. By making a suitable arrangement, an electric field lying in the plane of the figure, directed towards right with strength $\left(\dfrac{\pi}{10}\right)$MV/m is switched on exactly for an interval of 1 ms after every half second. Take the moment when the field is switched on for first time as $t=0$. Now, select the correct options.
\qfig{p2_phy_s2_q11_1}
\begin{qoptions}
\opt{A}{The maximum extension in the spring is 0.1414 m (nearly)}
\opt{B}{The spring never gets compressed.}
\opt{C}{The time taken by the block to start from mean position and reach the position of maximum elongation of spring is 0.75s.}
\opt{D}{The time taken by the block to start from mean position to the position of maximum elongation of spring is 0.5 s.}
\end{qoptions}
\qn{12} In a ``Positronium'' atom a pair of electron and positron orbit each other. If Bohr's atomic model is valid, then in comparison with Bohr's atomic hydrogen,\\
(assume that mass of proton $\ggg$ mass of electron)
\begin{qoptions}
\opt{A}{The Rydberg's constant for positronium is half}
\opt{B}{The ionisation energy for positronium is double}
\opt{C}{The radius of positronium is double}
\opt{D}{The kinetic energy of electron in the center of mass frame of positronium is one fourth that of the electron wrt proton in the H-atom}
\end{qoptions}
\sect{Section-II}
\qn{13} A circular disc of radius $R$ carries surface charge density $\sigma(r)=\sigma_0\left(1-\dfrac rR\right)$ where $\sigma_0$ is a constant and $r$ is the distance from the center of the disc. Electric flux through a large spherical surface that encloses the charged disc completely is $\phi_0$. Electric flux through another spherical surface of radius $\dfrac R4$ and concentric with the disc is $\phi$. Then the ratio $\dfrac{\phi_0}{\phi}$ is\underline{\hspace{1.2cm}}
\qn{14} A liquid at $30^\circ$C is poured very slowly into a Calorimeter that is at temperature of $110^\circ$C. The boiling temperature of the liquid is $80^\circ$C. It is found that the first 5 gm of the liquid completely evaporates. After pouring another 80 gm of the liquid the equilibrium temperature is found to be $50^\circ$C. The ratio of the Latent heat of the liquid to its specific heat will be \underline{\hspace{1.2cm}}$^\circ$C.\\
{[}Neglect the heat exchange with surrounding{]}
\qn{15} The pitch of a screw gauge is 1 mm and there are 50 divisions on its cap. When nothing is put in between the studs, 44$^{\text{th}}$ division of the circular scale coincides with the reference line and zero of main scale is not visible. When a plate is placed between the studs, the main scale reads 5 divisions and the circular scale reads 25 divisions. The true thickness of plate is $a\times10^{-5}$ m, find the value of $a$.
\qn{16} A source of sound and an observer are moving with constant velocities of 10 m/s and 60 m/s respectively. The source emits sound of frequency 657 Hz. At a certain moment, the distance between the source and observer is 415 m and the observer is moving in a direction making angle $37^\circ$ with direction of position vector of source with respect to observer while the source is headed directly towards the observer. If the frequency of sound (in Hz) perceived by the observer at an instant exactly 5 seconds after this moment is $x$, then find the value of $\dfrac{x}{225}$. The speed of sound is 225 m/s. $\cos37^\circ=4/5$.
\qn{17} Given that mass of Earth is $M$ and that of moon is $m$ and distance between them is $R$. Neglect effect of any other planet or sun. A small mass is projected from midway from Earth and moon its escape velocity is $P\sqrt{\dfrac GR(m+M)}$ find $P$.
\qn{18} A plane thick wall having uniform surface temperature along planes are $0^\circ$K and $T_0^\circ$K ($T_0=300$K) at $x=0$ and $x=x_0$ respectively. Thermal conductivity varies linearly with temperature $K=K_0(1+T)$. The temperature of wall at the plane $x=2x_0$ is $T'$ kelvin then find the value of $\dfrac{T'^2+2T'}{200}$
\partsec{2}{CHEMISTRY}
\sect{Section-I(i)}
\qn{1} In a CCP crystal of edge length `$a$' formed by an anion having radius $r_-$, assuming no distortion in the unit cell, the ratio of radius of tetrahedral void to octahedral void $\left(\dfrac{r_t}{r_o}\right)$ is
\begin{qoptions}
\opt{A}{$\dfrac{\sqrt6-2}{\sqrt8-2}$}
\opt{B}{$\dfrac{\sqrt3-2}{\sqrt8-2}$}
\opt{C}{$\dfrac{\sqrt6-2}{\sqrt2-2}$}
\opt{D}{$\dfrac{\sqrt3-2}{\sqrt2-2}$}
\end{qoptions}
\qn{2} The temperature coefficient of standard cadmium cell is $5.0\times10^{-5}\ \mathrm{VK^{-1}}$ whose emf at $27^\circ$C is 1.015 V. During the cell operation, the temperature will \underline{\hspace{0.8cm}} (The number of electrons in overall process involved are two).
\begin{qoptions}
\opt{A}{Increase}
\opt{B}{Decrease}
\opt{C}{Either}
\opt{D}{Remains constant}
\end{qoptions}
\qn{3} In the extraction of copper, the metal formed in the Bessemer's converter is due to the reaction :
\begin{qoptions}
\opt{A}{$\mathrm{Cu_2S\longrightarrow2Cu+S}$}
\opt{B}{$\mathrm{2Cu_2O\longrightarrow4Cu+O_2}$}
\opt{C}{$\mathrm{2Cu_2S+3O_2\longrightarrow2Cu_2O+2SO_2}$}
\opt{D}{$\mathrm{2Cu_2O+Cu_2S\longrightarrow6Cu+SO_2}$}
\end{qoptions}
\qn{4} For the carbohydrate given below, which is not a true statement :
\qfig{p2_chem_s1_q04_1}
\begin{qoptions}
\opt{A}{On reaction with Conc.$\mathrm{HNO_3}$, it gives an optically inactive compound}
\opt{B}{On reduction with $\mathrm{LiAlH_4}$, it gives a meso compound}
\opt{C}{Both L-Glucose and the above molecule will form same osazone}
\opt{D}{It is a $\mathrm{C_4}$ epimer of D (+) - Glucose}
\end{qoptions}
\sect{Section-I(ii)}
\qn{5} `Spin only' magnetic moment of Ni in $\mathrm{[Ni(dmg)_2]}$ is same as that found in
\begin{qoptions}
\opt{A}{Ni in $\mathrm{[NiCl_2(PPh_3)_2]}$}
\opt{B}{Mn in $\mathrm{[MnO_4]^-}$}
\opt{C}{Co in $\mathrm{[CoBr_4]^{2-}}$}
\opt{D}{Pt in $\mathrm{[Pt(H_2O)_2Br_2]}$}
\end{qoptions}
\qn{6} For the reaction
\qfig{p2_chem_s2_q06_1}
Where $K_1=6.93\times10^{-3}\ \mathrm{min^{-1}}$, $K_2=1.386\times10^{-2}\ \mathrm{min^{-1}}$\\
Initially reaction is started with 2 M of A. If both A, B, and C are optically active and have an optical rotation per unit concentration of $+40^\circ$, $-50^\circ$ and $+24^\circ$ respectively, then which of the following statement(s) is/are true
\begin{qoptions}
\opt{A}{The reaction has a half life of 45 min}
\opt{B}{When 50\% of A is consumed, the reaction mixture has an optical rotation of $+38.67^\circ$}
\opt{C}{When 75\% of the A is consumed, the reaction mixture has $0^\circ$ optical rotation}
\opt{D}{At infinite time, the reaction mixture is levo rotatory}
\end{qoptions}
\qn{7}
\qfig{p2_chem_s2_q07_2}
For the above sequence of reactions, which of the following statement(s) is/are true?
\begin{qoptions}
\opt{A}{X gives yellow precipitate with $\mathrm{I_2}$ in presence of NaOH}
\opt{B}{Y is reduced by Fehling's solution}
\opt{C}{Z gives silver mirror with $\mathrm{[Ag(NH_3)_2]^+}$}
\opt{D}{X reduces Tollen's reagent.}
\end{qoptions}
\qn{8}
\qfig{p2_chem_s2_q08_1}
Which of the following is/are not correct regarding above reaction ?
\begin{qoptions}
\opt{A}{If `G' is a +R group and present at para position, it will increase the rate of reaction}
\opt{B}{If `G' is a $-$R group and present at para position, it will increase the rate of reaction}
\opt{C}{If `G' is +I group and present at meta position, it will increase the rate of reaction}
\opt{D}{If `G' is $-$I group and present at meta position, it will increase the rate of reaction}
\end{qoptions}
\qn{9} In the given reaction sequence
\qfig{p2_chem_s2_q09_1}
Which of the following is correct?
\begin{qoptions}
\opt{A}{\qfigi{p2_chem_s2_q09_A}}
\opt{B}{\qfigi{p2_chem_s2_q09_B}}
\opt{C}{\qfigi{p2_chem_s2_q09_C}}
\opt{D}{\qfigi{p2_chem_s2_q09_D}}
\end{qoptions}
\qn{10} Which of the following is/are incorrect statement for $\mathrm{KMnO_4}$
\begin{qoptions}
\opt{A}{It can be prepared by the oxidation of $\mathrm{K_2MnO_4}$ by $\mathrm{H_2S_2O_8}$}
\opt{B}{On heating in acidic medium it disproportionates to give $\mathrm{MnO_2}$ and $\mathrm{MnO_4^{-2}}$}
\opt{C}{As an oxidising agent in neutral medium, $\mathrm{KMnO_4}$ equivalent mass is $\dfrac M3$}
\opt{D}{It can be used as primary standard solution in standardisation of oxalic acid sample in HCl medium.}
\end{qoptions}
\qn{11} For the reaction $\mathrm{P_4+SOCl_2\to A+B+SO_2}$, which of the following statements is/are false
\begin{qoptions}
\opt{A}{If A is phosphorous compound, then on hydrolysis it gives $\mathrm{H_3PO_4}$}
\opt{B}{If B is a compound of S, then is has similar structure that of $\mathrm{H_2O_2}$}
\opt{C}{In the above reaction, S undergoes disproportionation}
\opt{D}{A in solid state exists as $\mathrm{[PCl_4]^+[PCl_6]^-}$}
\end{qoptions}
\qn{12} The ionic conductance at infinite dilution of $\mathrm{Al^{+3}}$ and $\mathrm{SO_4^{-2}}$ ions are 60 and 80 $\Omega^{-1}\mathrm{cm^2eq^{-1}}$, respectively. The correct details regarding $\mathrm{Al_2(SO_4)_3}$ will be
\begin{qoptions}
\opt{A}{The molar conductance is 140 $\Omega^{-1}\mathrm{cm^2mol^{-1}}$}
\opt{B}{The equivalent conductance is 140 $\Omega^{-1}\mathrm{cm^2mol^{-1}}$}
\opt{C}{The molar conductance is 840 $\Omega^{-1}\mathrm{cm^2mol^{-1}}$}
\opt{D}{The molar conductance is 23.33 $\Omega^{-1}\mathrm{cm^2mol^{-1}}$}
\end{qoptions}
\sect{Section-II}
\qn{13} At $10^\circ$C, the osmotic pressure of urea solution is 500 mm. The solution is diluted and the temperature is raised to $25^\circ$C, when the osmotic pressure is found to be 105.3 mm. Then the extent of dilution is
\qn{14} A metal cube of edge length `$a$' cm is cut $n$ number of times. It is cut such that, per each cut, each cube is divided into 8 smaller cubes of same edge length. Then what should be the value of $n$, if the final surface area becomes 1048576 times of initial surface area
\qfullbegin
\qn{15} How many of the following molecules does not give positive haloform test.
\qfig{p2_chem_s3_q15_1}
\qn{16} How many of the following substrates will show rearrangement during $\mathrm{S_N1}$ reaction?
\qfig{p2_chem_s3_q16_1}
\qfullend
\qn{17} For the coagulation on 200 mL of $\mathrm{As_2S_3}$ solution, 10 mL of 1 M NaCl is required. What is the coagulating value of NaCl.
\qn{18} Number of possible isomers that can be formed (including stereo) for an octahedral complex containing $\mathrm{Co^{+3}}$ metal cation, two `en' ligands, two $\mathrm{NO_2^{-1}}$ and one $\mathrm{NO_3^{-1}}$ ions is $X$. Then $\dfrac X3$ will be
\partsec{3}{MATHEMATICS}
\sect{Section-I(i)}
\qn{1} \vspace{-4pt}{\small\[\lim_{n\to\infty}\sum_{r=1}^{7n}\left(\frac{r^6+6r^5+5r^4+4r^3+3r^2+2r+1}{7n^7+6n^6+5n^5+4n^4+3n^3+2n^2+7n+1}\right),\]}$k\in R$, $n\in N$ is
\begin{qoptions}
\opt{A}{$6^7$}
\opt{B}{$6^5$}
\opt{C}{$5^7$}
\opt{D}{$7^5$}
\end{qoptions}
\qn{2} If two points P \& Q on the hyperbola $9x^2-5y^2=1$ whose centre is C be such that CP is perpendicular to CQ, then $\dfrac{1}{CP^2}+\dfrac{1}{CQ^2}$ is
\begin{qoptions}
\opt{A}{Equal to the number of integers for which $n^4-20n^2+4$ is a prime number}
\opt{B}{Equal to the twice the number of integers for which $n^4-20n^2+4$ is a prime number}
\opt{C}{Equal to 4}
\opt{D}{Is equal to 2}
\end{qoptions}
\qn{3} The value of $S=\sec11^\circ\sec19^\circ-2\cot71^\circ$ equals
\begin{qoptions}
\opt{A}{$2\cot11^\circ$}
\opt{B}{$\tan19^\circ$}
\opt{C}{$2\tan11^\circ$}
\opt{D}{$\dfrac12\tan19^\circ$}
\end{qoptions}
\qn{4} Let $r_1,\ r_2,\ r_3,\ r_4$ be the four roots of the polynomial $x^4-4x^3+8x^2-7x+3$. Find the absolute value of $\dfrac{r_1^2}{r_2^2+r_3^2+r_4^2}+\dfrac{r_2^2}{r_1^2+r_3^2+r_4^2}+\dfrac{r_3^2}{r_1^2+r_2^2+r_4^2}+\dfrac{r_4^2}{r_1^2+r_2^2+r_3^2}$
\begin{qoptions}
\opt{A}{4}
\opt{B}{3}
\opt{C}{1}
\opt{D}{8}
\end{qoptions}
\sect{Section-I(ii)}
\qn{5} Let $ax^2+bx+a=0$ be the quadratic equation (where $a,b\in R$) has roots $\alpha,\beta$. Consider the family of lines $\left(b^2-2a^2+c^2\right)x+8a^2=2\left(b^2-c^2\right)y$ (where $c$ is a real number ). If $|\alpha-\beta|=2$ then
\begin{qoptions}
\opt{A}{The given family of lines always passes through a fixed point $\left(\dfrac{-4}{7},\dfrac27\right)$}
\opt{B}{The given family of lines always passes through a fixed point $(2\,,-4)$}
\opt{C}{Number of lines belongs to the given family which makes a triangle with coordinate axes of area 16 are only 4}
\opt{D}{Number of lines belongs to the given family which makes a triangle with coordinate axes of area 2 are only 2}
\end{qoptions}
\qn{6} The vector $\hat i+x\hat j+3\hat k$ is rotated about its initial point through an angle of $\cos^{-1}\!\left(\dfrac{11}{14}\right)$ and it's magnitude is doubled. If the vector in the new position is given by $4\hat i+(4x-2)\hat j+2\hat k$ then which of the following can't be the value of $x$ ?
\begin{qoptions}
\opt{A}{$-\dfrac23$}
\opt{B}{$\dfrac23$}
\opt{C}{$-\dfrac{20}{17}$}
\opt{D}{2}
\end{qoptions}
\qn{7} Let $f(x)$ is a differentiable function such that $u,v\in R$\\
$2u\,f\!\left(u+v^2\right)+v^3+u^4v=2v\,f\!\left(u^2+v\right)+u^3+uv^4$\\
Consider the curve $P$ defined by\\
$P:\ f(x+y)-x-y+\dfrac12=2f\!\left(\sqrt{xy}\right)$, then -
\begin{qoptions}
\opt{A}{Area bounded by curve $P$ with coordinate axis is $\dfrac{4-\pi}{4}$}
\opt{B}{Range of the function $f(x)$ is $R^+$}
\opt{C}{Number of normals which can be drawn from point $(1,1)$ to the curve $P$ will be 4}
\opt{D}{Length of tangent drawn from origin to the curve $P$ is 1}
\end{qoptions}
\qn{8} Let $I=\displaystyle\int_0^1\sqrt{\dfrac{1+\sqrt x}{1-\sqrt x}}\,dx$ and $J=\displaystyle\int_0^1\sqrt{\dfrac{1-\sqrt x}{1+\sqrt x}}\,dx$ then which statement(s) are correct
\begin{qoptions}
\opt{A}{$I+J=2$}
\opt{B}{$I-J=\pi$}
\opt{C}{$J=\dfrac{2+\pi}{2}$}
\opt{D}{$J=\dfrac{4-\pi}{2}$}
\end{qoptions}
\qn{9} If A, B, C are the angles of an acute angled triangle ABC and
{\footnotesize\[D=\begin{vmatrix}(\tan B+\tan C)^2&\tan^2A&\tan^2A\\ \tan^2B&(\tan A+\tan C)^2&\tan^2B\\ \tan^2C&\tan^2C&(\tan A+\tan B)^2\end{vmatrix},\]}
then the value of $\dfrac{D}{1000}$ can be
\begin{qoptions}
\opt{A}{$-1$}
\opt{B}{0}
\opt{C}{2}
\opt{D}{3}
\end{qoptions}
\qn{10} An electrical apparatus with three components $C_1$, $C_2$ and $C_3$ as in the diagram. The components works independently of each other with probabilities $\dfrac{8}{10},\dfrac{9}{10}$ and $\dfrac34$. For apparatus to be operational current flows from left to right. Let $X$ denote the event that apparatus is operational and $X_1,X_2,X_3$ denotes respectively the events that the component $C_1$, $C_2$ and $C_3$ are functioning. Which of the following is(are) true ?
\qfig{p2_math_s2_q10_1}
\begin{qoptions}
\opt{A}{$P$(exactly two components are functioning $|X)=1$}
\opt{B}{If $P(X)$ is $\dfrac ab$ ($a,b$ are co-prime) then $b-a$ is 7}
\opt{C}{$P(X|X_3)=1$}
\opt{D}{$P(X|X_2)=19/20$}
\end{qoptions}
\qn{11} Point A lies on x-axis and point B lies on $y=x$. If $PA.PB$ is minimum where $P$ divides segment AB internally, then (given, $P$ is $\left(1,2-\sqrt3\right)$ and `O' is origin), -
\begin{qoptions}
\opt{A}{$OA=2\sqrt2-\sqrt6+\sqrt3-1$}
\opt{B}{$OB=2\sqrt2-\sqrt6+\sqrt3-1$}
\opt{C}{$\angle OAB=45^\circ$}
\opt{D}{$\angle OBA=22.5^\circ$}
\end{qoptions}
\qn{12} For the function $f(x)=\sqrt{4x^2-4x+2}+|x|$ which of the following holds goods
\begin{qoptions}
\opt{A}{The least value of the function is $\dfrac{\sqrt3+1}{2\sqrt2}$}
\opt{B}{The least value of function is $\dfrac{\sqrt3+1}{\sqrt2}$}
\opt{C}{The function takes its least value at $x=\dfrac{3-\sqrt3}{6}$}
\opt{D}{Function has exactly one point of minima}
\end{qoptions}
\sect{Section-II}
\qn{13} Let the function $g(x)$ where $g:R\to R$, be continuous in $[0,\infty)$ and positive in $(0,\infty)$ satisfying the condition $g(1)=1$ and $\dfrac12\displaystyle\int_0^xg^2(t)\,dt=\dfrac1x\left(\displaystyle\int_0^xg(t)\,dt\right)^2$ then the value of $\left(g(4)\right)^{\sqrt2-1}$ is
\qn{14} If sum and the difference of the roots of the characteristic equation of a square matrix of order $2\times2$ are 4 and 2 respectively then the determinant of the matrix will be
\qn{15} If $\begin{bmatrix}4a^2&4a&1\\4b^2&4b&1\\4c^2&4c&1\end{bmatrix}\begin{bmatrix}f(-1)\\f(1)\\f(2)\end{bmatrix}=\begin{bmatrix}3a^2+3a\\3b^2+3b\\3c^2+3c\end{bmatrix}$, $f(x)$ is a quadratic function and its maximum value occurs at a point V. A is a point of intersection of $y=f(x)$ with x-axis and point B is such that chord AB subtends a right angle at V. Let A be the area enclosed by $f(x)$ and chord AB. Find $[A]$ (where $[\cdot]$ denotes the greatest integer function)
\qn{16} The value of the series $\dfrac{\displaystyle\sum_{r=0}^{100}(-1)^r\,{}^{100}C_r\,(r)^{50}}{(50)^{100}}$ is equal to
\qn{17} Figure shows a tetrahedron $OA_1A_2A_3$. Altitudes $OP$ and $A_1Q$ are drawn from $O$ and $A_1$ on opposite triangular faces. It is given that $OP$ and $A_1Q$ intersect at $H$. Let the lengths of $OA_1$, $A_2A_3$ be $l_1$ and $l_2$ respectively. If shortest distance between $OA_1$ and $A_2A_3$ is $d$ and volume of tetrahedron is $V$. Value of $\dfrac{d\ell_1\ell_2}{V}$ is equal to
\qfig{p2_math_s3_q17_1}
\qn{18} Let $S_n=\displaystyle\sum_{r=1}^{n}r.2^{\frac r2\left(1+(-1)^r\right)}.3^{\frac r2\left(1-(-1)^r\right)}$. If $S_{20}=\alpha.3^{21}+\beta.2^{22}+\dfrac{391}{288}$, then the value of $\left[\dfrac{32\alpha-9\beta}{4}\right]$ is (where $[.]$ denotes the greatest integer function)
\end{multicols}
\end{document}
